%This is an example .tex document for the Warner style file, adapted from the Lankton theme
% Zach Warner, University of Wisconsin-Madison
% 12 December 2013
% contact: zwarner@wisc.edu

% NB: I prefer the customization available by putting the style elements directly in the .tex document, but this will work for Wisconsin people looking for a quick poster.

\documentclass[serif,mathserif,final]{beamer}
\usetheme{Warner}
\usepackage{amsmath,amsfonts,amssymb,pxfonts,eulervm,xspace}
\usepackage{graphicx,tikz}
\usepackage{wrapfig}
\usepackage{textpos} 
\usepackage{ragged2e}
\usepackage[orientation=landscape,size=custom,width=48,height=36,scale=.6,debug]{beamerposter}

%note that to change the logo in the header, you have to edit the beamerthemeWarner.sty file.


%-- Header and footer information ----------------------------------
\newcommand{\footleft}{PolMeth Annual Conference 2013}
\newcommand{\footright}{zwarner@wisc.edu}
\newcommand{\footmid}{\today}
\title{A Wisconsin Poster Template}
\author{Zach Warner}
\date{\today}
\institute{zwarner@wisc.edu} % well, that's a lie. But email makes more sense here
%-------------------------------------------------------------------


%-- Main Document --------------------------------------------------
\begin{document}
\begin{frame}{}
  \begin{columns}[t]

    %-- Column 1 ---------------------------------------------------
    \begin{column}{0.32\linewidth}

      %-- Block 1-1
      \begin{block}{Abstract}
        Some text. Don't worry, the white space above the first texts blocks will disappear as the columns fill up.
      \end{block}

      %-- Block 1-2
       \begin{block}{Motivation}      
       A great place for some ponderous questions, and perhaps a model or two. If using a model, it might get shortened due to the columns. I'd use \texttt{multiline} to make it work. 
      \end{block}



     %-- Block 1-3
     \begin{block}{A Test}
         Run it.
         \begin{itemize}
          \item<circ@1->
          This code makes sure the colors of the circles match the color theme.
           \item<circ@1->
           If you change this to fit your school, you can set \texttt{fg=black} to whatever rgb coordinates you desire.
            \item<circ@1->
            While I'm in this list, note that the normal line break of won't work here. 
			\item<circ@1->            
            Instead, use \texttt{par}, a tilde, and another \texttt{par}.
         \end{itemize}
         One final snafu: the crest won't be where it should be whenever you compile for the first time. Always compile twice when you make adjustments to it.
     \end{block}
     
     
    \end{column}%1

    %-- Column 2 ---------------------------------------------------
    \begin{column}{0.32\linewidth}

      %-- Block 2-2
      \begin{block}{Results}
      Put your main results, particularly attractive visual displays, here.
       \begin{figure}[htb]
 % \vspace{-5pt} usually a figure will leave too much white space, so use negative vspace to adjust it to your taste
  \centering
          \includegraphics[width=\columnwidth]{correlation.pdf}
        \end{figure}
        Source: \texttt{http://xkcd.com/552/}.
      \end{block}

  %-- Block 2-3
      \begin{block}{Robustness}
They don't think it be like it is, but it do.
      \end{block}    

    \end{column}%2

    %-- Column 3 ---------------------------------------------------
    \begin{column}{0.32\linewidth}

      %-- Block 3-1
      \begin{block}{Extensions 1}
  This is a great place to insert tables and graphics side-by-side using \texttt{column} or \texttt{wrapfig}. Make sure you use \texttt{booktabs} for tables to make them prettier.
      \end{block}

      %-- Block 3-2
      \begin{block}{Extensions 2}
      One last figure, possibly just to show some interesting avenues for future research. 
      \end{block}

	%-- Block 3-3
	 \begin{block}{Conclusion}
	Based on the template ``Lankton'' by Shawn Lankton, available at \texttt{http://		www.shawnlankton.com/2008/06/latex-beamer- \\ poster-theme-and-template/}. \par ~ \par
Feel free to email me questions, comments, hatemail, or thanks, if so inclined.
      \end{block}

    \end{column}

  \end{columns}
  
%--- The crest----------------------------------------------------
  \begin{tikzpicture}[remember picture,overlay] 
            \node[opacity=.1, at=(current page.south west),anchor=south west,inner sep=-160pt] {
                \includegraphics[width=.5\paperheight]{uwcrest.pdf}
            };
        \end{tikzpicture}
        % higher numbers for opacity mean a brighter picture, but this can make text harder to read
        % the anchor is the bottom left column, inner sep moves it so it bleeds off the page.
        % adjust inner sep according to taste
        
        %note that the first time you compile, the crest won't be where you want it to be. Compile a second time and it should be fine.
        
\end{frame}
\end{document}
